%%%%%%%%%%%%%%%%% Allgemeines Dokumentenformat %%%%%%%%%%%%%%%%%

\documentclass[a4paper,12pt,headsepline]{scrartcl}

%%%%%%%%%%%%%%%%% Einbinden von Paketen %%%%%%%%%%%%%%%%%

% Grafiken aus PNG Dateien einbinden
\usepackage{graphicx}

% Deutsche Sonderzeichen benutzen 
\usepackage{ngerman}

% deutsche Silbentrennung
\usepackage[ngerman]{babel}

% Eurozeichen einbinden
\usepackage[right]{eurosym}

% Umlaute unter UTF8 nutzen
\usepackage[utf8]{inputenc}

% Zeichenencoding
\usepackage[T1]{fontenc}

% Schriftart
\usepackage{lmodern}
\usepackage{fix-cm}

% mehrseitige Tabellen ermöglichen
\usepackage{longtable}

% Packet für Seitenrandabständex und Einstellung für Seitenränder
\usepackage{geometry}

% Paket für Boxen im Text
\usepackage{fancybox}

% bricht lange URLs "schön" um
\usepackage[hyphens,obeyspaces,spaces]{url}

% Paket für Textfarben
\usepackage{color}

% Mathematische Symbole importieren
\usepackage{amssymb}

% erzeugt Inhaltsverzeichnis mit Querverweisen zu den Kapiteln (PDF Version)
\usepackage[bookmarksnumbered,hyperfootnotes=false]{hyperref}

% neue Kopfzeilen mit fancypaket
\usepackage{fancyhdr} %Paket laden

% für Tabellen
\usepackage{array}

% Paket für Zeilenabstand
\usepackage{setspace}

% für Bildbezeichner
\usepackage{capt-of}

% für Stichwortverzeichnis
\usepackage{makeidx}

% für Listings
\usepackage{listings}

% Glossar und Akronyme
\usepackage[acronym, toc]{glossaries}

%%%%%%%%%%%%%%%%% Seitenformatierung %%%%%%%%%%%%%%%%%

% Layout
\geometry{left=3.5cm, right=2cm, top=2.5cm, bottom=2cm}

% Kopfzeile
\pagestyle{fancy} %eigener Seitenstil
\fancyhf{} %alle Kopf- und Fußzeilenfelder bereinigen
\fancyhead[L]{\nouppercase{\leftmark}} %Kopfzeile links
\fancyhead[C]{} %zentrierte Kopfzeile
\fancyhead[R]{\thepage} %Kopfzeile rechts
\renewcommand{\headrulewidth}{0.4pt} %obere Trennlinie

% Festlegung Art der Zitierung - Havardmethode: Abkuerzung Autor + Jahr
\bibliographystyle{alphadin}

% Schaltet den zusätzlichen Zwischenraum ab, den LaTeX normalerweise nach einem Satzzeichen einfügt.
\frenchspacing

% Formatierung für Listings
\lstset{numbers=left, numberstyle=\tiny, numbersep=5pt, keywordstyle=\color{black}\bfseries, stringstyle=\ttfamily,showstringspaces=false,basicstyle=\footnotesize,captionpos=b}
\lstset{language=java}

% Indexerstellung
\makeindex

% Glossar erstellen
\makeglossaries

% Meta-Informationen erstellen 
\hypersetup{
  pdftitle    = {Titel},
  pdfsubject  = {Um was geht es},
  pdfauthor   = {Autor bzw. Autoren},
  pdfkeywords = {Stichwort1, Stichwort2, ...}
}

% Disable single lines at the start of a paragraph (Schusterjungen)
\clubpenalty = 10000
% Disable single lines at the end of a paragraph (Hurenkinder)
\widowpenalty = 10000
\displaywidowpenalty = 10000

%%%%%%%%%%%%%%%%% Definitionen %%%%%%%%%%%%%%%%%

% Glossar-Definitionen
\include{definitionen/glossar}

% Variablen-Definitionen
\include{definitionen/variablen}

% Silbentrennungs-Definitionen
\include{definitionen/trennung}

%%%%%%%%%%%%%%%%% Inhalt %%%%%%%%%%%%%%%%%

\begin{document}

% Titelseite %
\include{extras/deckblatt}

% Sperrvermerk
\include{extras/sperrvermerk}

% Eidesstattliche Erklärung
\include{extras/erklaerung}

% Einleitung / Abstract
\include{extras/abstract}

% Inhaltsverzeichnis
\newpage
\tableofcontents

\newpage
\fancyhead[L]{Abbildungs- / Tabellen- / Listingverzeichnis}
% Abbildungsverzeichnis
\addcontentsline{toc}{section}{Abbildungsverzeichnis} % Soll im Inhaltsverzeichnis auftauchen
\listoffigures
% Tabellenverzeichnis
\addcontentsline{toc}{section}{Tabellenverzeichnis}
\listoftables
% Listingverzeichnis
\addcontentsline{toc}{section}{Listingverzeichnis}
\renewcommand{\lstlistlistingname}{Listingverzeichnis} % eigenen Listingverzeichnis-Titel setzen
\lstlistoflistings

% Glossar
\newpage
\fancyhead[L]{Glossar} %Kopfzeile links
\printglossary

% Abkürzungsverzeichnis
\newpage
\fancyhead[L]{Abkürzungsverzeichnis}
\printglossary[type=\acronymtype, title={Abkürzungsverzeichnis}]

% einzelne Kapitel
\newpage
\fancyhead[L]{\nouppercase{\leftmark}} %Sectiontitle in die Kopfzeile
\include{kapitel/1_abbildungen}
\include{kapitel/2_tabellen}
\include{kapitel/3_quellcode}
\include{kapitel/4_formeln}
\include{kapitel/5_zitate}
\include{kapitel/6_begriffe}

% Literaturverzeichnis
\newpage
\addcontentsline{toc}{section}{Literaturverzeichnis}
\bibliography{definitionen/literatur}

% Anhang
\include{extras/anhang}

\end{document}